\section*{Chapitre \ref{ch:rs:risk-models} --- Modèles de risque}

\begin{solution}{exo:rs:risk-models:1}
$\sqrt{3.26^2 + (1.23 \times 7.2)^2} = 9.4\%$, contre une volatilité réalisée de 9.2~\%~; la \omterm{def:rs:risk-models:bias}{statistique de biais} vaut environ $9.4/3.26 = 2.9$ (la valeur mesurée est 2.86).
\end{solution}

\begin{solution}{exo:rs:risk-models:2}
$1 \pm \sqrt{2/60} = 1 \pm 0.18$~; $1 \pm \sqrt{2/2\,016} = 1 \pm 0.031$.
\end{solution}

\begin{solution}{exo:rs:risk-models:3}
Les indicatrices sectorielles somment à l'exposition au facteur pays (une pour chaque action), si bien que la régression est singulière~: une constante peut passer du facteur pays
aux secteurs. Imposer que les rendements sectoriels pondérés par la capitalisation soient nuls fait du facteur pays le marché pondéré par la capitalisation, et de chaque \omterm{def:rs:risk-models:style}{facteur sectoriel}
le rendement du secteur net du marché (un portefeuille sectoriel neutre en dollars).
\end{solution}

\begin{solution}{exo:rs:risk-models:4}
Chaque prévision est la variance vraie plus un bruit d'estimation. Trier selon la prévision revient en partie à trier selon ce bruit~: le premier décile regroupe les actions dont les résidus récents
ont été par chance grands (variance vraie inférieure à la prévision), le dernier décile regroupe celles dont ils ont été par chance faibles. La variance réalisée revient vers la vraie, si bien que le premier décile est surestimé
(0.93), et le dernier décile sous-estimé (1.13). Rapprocher les prévisions d'une moyenne de groupe défait une partie du tri.
\end{solution}

\begin{solution}{exo:rs:risk-models:5}
L'ajustement multiplie la covariance des facteurs par le biais transversal récent de ces facteurs. Le risque du portefeuille de marché est presque entièrement un risque factoriel, dont les régimes
l'ajustement suit. Le risque des portefeuilles aléatoires est surtout spécifique~: ils n'y gagnent rien et héritent du bruit du multiplicateur sur leur petite part factorielle.
\end{solution}

\begin{solution}{exo:rs:risk-models:6}
$\sum_k x_k (Fx)_k = x^\top F x$. Une contribution est négative lorsque l'exposition à un facteur couvre les autres~: $x_k$ et $(Fx)_k$ sont de signes opposés.
\end{solution}

\begin{solution}{exo:rs:risk-models:7}
Exécuter \texttt{rs\_riskmodel.summary(0.3)}. L'exposition au momentum baisse, si bien que le modèle sans momentum sous-estime moins~: \omterm{def:rs:risk-models:bias}{statistiques de biais} de 1.98
sans momentum, 1.02 avec tous les facteurs et 1.48 pour le modèle statistique (volatilité réalisée de 6.53~\%, prévision de 3.34~\% sans momentum).
\end{solution}

\begin{solution}{exo:rs:risk-models:8}
Le $R^2$ porte sur des actions isolées, dont la variance quotidienne est surtout spécifique. Le risque d'un portefeuille est la part factorielle, qui ne se diversifie pas, plus une part spécifique qui
diminue avec le nombre de titres~; la valeur du modèle tient à la prévision des deux, mesurée par les \omterm{def:rs:risk-models:bias}{statistiques de biais} (1.04 sur le portefeuille exposé, 1.00 en moyenne pour les portefeuilles aléatoires),
et non par l'ajustement de la régression.
\end{solution}

\begin{solution}{pb:rs:risk-models:1}
\begin{enumerate}
\item Facteur pays (un pour chaque action)~; dix \omterm{def:rs:risk-models:style}{facteurs sectoriels} (appartenance)~; bêta (régression sur un an glissant contre le marché pondéré par la capitalisation, mensuelle)~; taille, valeur, momentum (taille figée à la cotation~; valeur issue du dernier bilan déposé~; momentum issu du rendement 12--1 mois~; toutes connues à la clôture précédente).
\item La variance spécifique décroît avec la capitalisation~; les pondérations approchent les moindres carrés généralisés et donnent au facteur pays le sens du marché pondéré par la capitalisation.
\item 17.8~\%~; entre 8.0~\% et 8.7~\%~; 7.2~\%.
\item 0.19~: la variance quotidienne d'une action est surtout spécifique, si bien que le \omterm{def:rs:risk-models:specific}{risque spécifique} compte même dans un portefeuille de deux cents titres.
\item Il multiplie la covariance factorielle par une moyenne pondérée exponentiellement du biais transversal~; les \omterm{def:rs:risk-models:bias}{statistiques de biais} glissantes sur un an du marché couvrent
0.56--1.32 au lieu de 0.50--1.55, et 44~\% des fenêtres tombent hors de la bande au lieu de 66~\%.
\item Leur risque est surtout spécifique~; la part des portefeuilles aléatoires dans la bande est de 92~\% avec lui et de 98~\% sans lui.
\item De 1.13 dans le décile le plus bas à 0.93 dans le plus haut~; un rapprochement de 10~\% vers la moyenne du décile de taille donne de 0.97 à 1.04.
\item Les chocs de Student-$t$ et les sauts des jours d'annonce de résultats de la simulation, que la variance pondérée exponentiellement intègre avec retard.
\item Un score égal à 0.5 fois le score de momentum plus $\sqrt{0.75}$ fois du bruit~; les cent titres en tête et en queue~; neutralisé par rapport au modèle sans momentum~;
exposition brute de 2~; reconstruit chaque mois.
\item Tous les facteurs~: prévision de 9.23~\%, volatilité réalisée de 9.20~\%, \omterm{def:rs:risk-models:bias}{statistique de biais} 1.04. Sans momentum~: 3.26~\%, 9.20~\%, 2.86. Modèle statistique~: 5.59~\%, 9.20~\%, 1.68.
\item 76~\% de la variance prévue au jour 1\,500, pour une exposition de 1.03 (racine de la moyenne quadratique 1.23 sur les huit ans).
\item Le momentum est une direction de variation commune, que les composantes principales retrouvent~; mais ses chargements sont réinitialisés chaque mois, alors que les composantes calculées sur une année glissante
supposent des chargements qui restent fixes.
\item \textbf{Résultat nommé.} \omterm{def:rs:risk-models:bias}{Statistique de biais} 2.86 pour le modèle sans momentum~; volatilité réalisée égale à 2.83 fois la prévision (9.20~\% contre 3.26~\%).
\item En estimant un modèle statistique à côté du modèle fondamental (5.59~\% contre 3.26~\%), ou en régressant les rendements du portefeuille sur des \omterm{def:rs:risk-models:families}{rendements des facteurs} candidats.
\item \omterm{def:rs:risk-models:bias}{Statistiques de biais} des portefeuilles aléatoires, des portefeuilles factoriels et des portefeuilles de la firme, glissantes et sur l'ensemble de l'échantillon, comparées au modèle en place.
\item La bande est trop étroite~: des prévisions correctes sortent de cette bande plus de 5~\% du temps (86~\% à l'intérieur pour un kurtosis de 5 sur 120 périodes, selon les chiffres de MSCI).
\item Un modèle fondamental auquel on ajoute des facteurs statistiques (par exemple les composantes principales de ses résidus)~; cela en vaut la peine lorsque les portefeuilles prennent des risques que les caractéristiques
ne décrivent pas.
\item Les portefeuilles aléatoires font la moyenne sur les expositions~; un modèle peut avoir raison en moyenne et se tromper pour l'exposition particulière que concentre un portefeuille réel.
\item Le portefeuille que détient la firme, et les portefeuilles construits à partir de ses signaux, car un signal peut s'exposer à un facteur absent du modèle.
\item Les risques auxquels il a donné un nom, et rien d'autre.
\end{enumerate}
\end{solution}

\begin{solution}{iq:rs:risk-models:1}
Fondamental~: expositions observées, \omterm{def:rs:risk-models:families}{rendements des facteurs} obtenus par régression transversale~; interprétable, stable, utile pour l'attribution et les contraintes. Statistique~: les deux
sont estimés à partir des rendements~; il repère une variation commune non nommée et s'adapte, mais ses facteurs sont difficiles à interpréter et instables. Privilégier le modèle fondamental pour la construction et le reporting,
le modèle statistique servant de contrôle et aux horizons courts~; les hybrides combinent les deux.
\end{solution}

\begin{solution}{iq:rs:risk-models:2}
Chaque jour, régression par moindres carrés pondérés des rendements transversaux sur les expositions de la clôture précédente, la pondération étant la racine carrée de la capitalisation, et
les rendements sectoriels pondérés par la capitalisation étant contraints à zéro, car les indicatrices pays et sectorielles sont colinéaires~; la contrainte fait du facteur pays le marché et
des \omterm{def:rs:risk-models:style}{facteurs sectoriels} les rendements sectoriels nets de celui-ci.
\end{solution}

\begin{solution}{iq:rs:risk-models:3}
Décomposer le P\&L réalisé sur les facteurs du modèle et le résidu~; chercher une composante commune dans les résidus (composantes principales des résidus des titres du portefeuille
ou régression sur des rendements de facteurs candidats)~; comparer à la prévision d'un modèle statistique~; vérifier le régime (un changement de volatilité que le modèle rattrape avec retard).
\end{solution}

\begin{solution}{iq:rs:risk-models:4}
Standardiser les rendements de nombreux portefeuilles par leurs prévisions et calculer l'écart-type (la \omterm{def:rs:risk-models:bias}{statistique de biais}), globalement et sur des fenêtres glissantes, avec sa bande
$1 \pm \sqrt{2/T}$. Pièges~: les queues épaisses élargissent la vraie bande~; un seul portefeuille constitue une preuve faible~; la moyenne sur les portefeuilles masque les erreurs sur des expositions particulières.
\end{solution}

\begin{solution}{iq:rs:risk-models:5}
L'optimiseur cherche les directions que la covariance estimée juge les moins risquées, et ce sont de façon disproportionnée celles où l'erreur d'échantillonnage a rendu l'estimation trop
faible~; Menchero, Wang et Orr corrigent les biais des portefeuilles propres.
\end{solution}

\begin{solution}{iq:rs:risk-models:6}
Entrées~: prix et rendements disponibles à chaque date, capitalisations, classifications sectorielles et descripteurs. Ordre~: expositions à la clôture~; régression factorielle du jour et
résidus~; covariance, multiplicateur de régime, variances spécifiques~; prévisions. Contrôles~: contrainte respectée, $R^2$ et \omterm{def:rs:risk-models:families}{rendements des facteurs} dans leurs plages, pas de saut d'exposition
sans cause, \omterm{def:rs:risk-models:bias}{statistiques de biais} suivies. Conserver chaque jour, en versions successives, les expositions, les \omterm{def:rs:risk-models:families}{rendements des facteurs}, les résidus et les prévisions, afin qu'une prévision passée puisse être reproduite.
\end{solution}
